% Generated draft: requires mathematical and source review.
\documentclass[11pt]{article}
\usepackage[margin=25mm]{geometry}
\usepackage{fontspec}
\setmainfont{Latin Modern Roman}
\usepackage{amsmath,amssymb,mathtools,unicode-math}
\setmathfont{Latin Modern Math}
\usepackage{enumitem,needspace,xurl,hyperref}
\setlength{\parindent}{0pt}
\setlength{\parskip}{5pt}
\setlength{\emergencystretch}{3em}
\newcommand{\N}{\mathbb{N}}
\newcommand{\Z}{\mathbb{Z}}
\newcommand{\Q}{\mathbb{Q}}
\newcommand{\R}{\mathbb{R}}
\newcommand{\C}{\mathbb{C}}
\newcommand{\F}{\mathbb{F}}
\newcommand{\A}{\mathbb{A}}
\newcommand{\PP}{\mathbb{P}}
\newcommand{\E}{\mathbb{E}}
\newcommand{\eps}{\varepsilon}
\newcommand{\dd}{\,\mathrm{d}}
\begin{document}
\section*{UnsolvedMath ID: 30003712}
\textbf{Source problem code:} catalogue entry
\subsection{Definitions and mathematical statement}
The supplied record asks three related questions but omits the ambient space of the coordinates, the range of the parameters, and the definition of \(M_{d,n}\). Consequently, the following complex interpretation of the first question is conditional and requires source confirmation.

For an integer \(n\geq 2\), put
\[
Q_n(x_1,\ldots,x_n)=\prod_{1\leq i<j\leq n}(x_i-x_j)^2,
\qquad
F_n=\{x\in\C^n:Q_n(x)=1\},
\]
and let
\[
\mu_{n(n-1)}=\{\zeta\in\C^\times:\zeta^{n(n-1)}=1\}.
\]
Here \(S_n\) denotes the symmetric group on \(\{1,\ldots,n\}\). Coordinate permutation and simultaneous scalar multiplication give commuting actions, with the convention
\[
((\sigma,\zeta)\cdot x)_j=\zeta x_{\sigma^{-1}(j)}.
\]
These preserve \(F_n\), since \(Q_n\) is symmetric and homogeneous of degree \(n(n-1)\). Interpret \(H^i(F_n,\Q)\) as singular cohomology for the complex topology, with the left group action \(g\cdot\alpha=(g^{-1})^*\alpha\). These ambient-space, topology, parameter-range, and action conventions are explicit interpretations; the supplied statement does not specify them.

\begin{enumerate}
\item Under this interpretation, determine, for every integer \(n\geq2\), the graded rational representation
\[
H^*(F_n,\Q)=\bigoplus_{i\geq0}H^i(F_n,\Q)
\]
of \(S_n\times\mu_{n(n-1)}\).

\item Determine whether a category analogous to \(\mathrm{FI}\) makes the sequence \(\{H^i(F_n,\Q)\}_n\), for each fixed cohomological degree \(i\), a finitely generated module. Here \(\mathrm{FI}\) is the category of finite sets and injections. For a specified category \(\mathcal C\), a rational \(\mathcal C\)-module means a covariant functor \(V:\mathcal C\to\mathrm{Vect}_{\Q}\); it is finitely generated if finitely many vectors in its values generate all values under linear combinations and the maps assigned to morphisms. To make this question precise, one must specify the intended category, its objects corresponding to the parameters \(n\), the functorial maps on cohomology, and how these maps relate to the displayed group actions. The phrase ``analogous to \(\mathrm{FI}\)'' supplies no criterion selecting an admissible category. The intended range of \(i\) and any uniformity requirement also need confirmation.

\item Determine whether \(M_{d,n}\) has a natural compactification. The supplied record does not define \(M_{d,n}\), the parameters \(d,n\), their ranges, the relevant geometric category, or the meaning of ``natural.'' This part therefore remains an unresolved formulation, rather than a precise compactification problem.
\end{enumerate}

The supplied dated curation labels the problem partially solved and attributes representation stability and stable integral homology results for pure-braid Milnor fibers to Miller and Tosteson, describing these as answering the finite-generation direction. It describes the full product-group module structure and compactification questions as remaining open. These are supplied literature assessments, not independently verified claims; their exact scope and applicability to the displayed squared-product fiber require review.
\subsection{Short English statement}
Determine the rational cohomology representations of the displayed braid Milnor fibers, identify the intended category governing finite generation across \(n\), and formulate the proposed compactification question for \(M_{d,n}\). The complex ambient space and natural group actions are conditional interpretations. The category and \(M_{d,n}\) are unspecified, and the supplied partially solved status requires independent review.
\subsection{Sources}
\begin{itemize}
\item Local UnsolvedMath registry, permanent ID \(30003712\), exact source problem code \texttt{catalogue entry}, title ``Representation Stability and Compactifications for Braid Milnor Fibers,'' in the collection ``Oberwolfach Reports Open Problems.'' The supplied record lists proposed year \(2018\) and names no individual proposer.
\item Supplied report citation: \emph{Topology of Arrangements and Representation Stability (2019)}, Oberwolfach Reports, DOI \texttt{10.4171/owr/2018/2}, pp.\ 43--123. \url{https://doi.org/10.4171/owr/2018/2}. The literature-source label is ``Oberwolfach Report 2/2018.'' The citation's \(2019\) designation and the report's \(2018\) designation are retained without reconciliation. The background describes Oberwolfach Reports as documenting mathematical workshops and their open-problem discussions.
\item Supplied literature source: Miller--Tosteson, \emph{Transactions of the AMS}. \url{https://doi.org/10.1090/tran/8466}. No article title or theorem identifiers are supplied.
\item Literature assessment checked at \texttt{2026-08-21}: supplied dated curation, not independently verified. It labels the problem ``partially solved,'' attributes representation stability and stable integral homology for pure-braid Milnor fibers to Miller and Tosteson, and says that this answers the finite-generation direction while the full \(S_n\) times root-of-unity module structure and compactification questions remain open. The background expressly calls this a dated literature triage derived from supplied OWR curation and requests independent verification of time-sensitive status and cited claims.
\end{itemize}
\end{document}
